\chapter{Fourier Computations and the Basel Sum}\label{ch:fourier}

We need not begin with a space of arbitrary functions or a completeness
statement about an orthogonal basis.  A finite trigonometric sum is a
computation, and a positive finite kernel gives an explicit approximation
bound.  One worked periodic function will then identify the reciprocal-square
sum.

\section{Finite orthogonality}
Write $e_k(x)=E(i\pi kx)=C(kx)+iS(kx)$ for integer $k$.  The trigonometric
FTC gives
\[
 \frac12\int_{-1}^1 e_k(x)\,dx
 =\begin{cases}1,&k=0,\\0,&k\ne0.\end{cases}
\]
Products satisfy $e_ke_l=e_{k+l}$.  These identities determine integrals
of every finite trigonometric polynomial by finite linear algebra.

For a supplied period-$2$ continuous computation $f$, with a uniform
continuity certificate on the periodic interval, define
\[
 \widehat f(k)=\frac12\int_{-1}^1f(x)e_{-k}(x)\,dx.
\]
Each coefficient is an ordinary integral computation.  The word periodic
includes matching endpoint values and the continuity bound across that
join; this data is checked for the particular functions used here.

\section{A positive kernel with an explicit tail}
For $N\ge1$, define
\[
 K_N(t)=\frac1N\left|\sum_{j=0}^{N-1}e_j(t)\right|^2
       =\sum_{|k|<N}\left(1-\frac{|k|}{N}\right)e_k(t).
\]
The identity follows by counting pairs with a given difference.  Thus
$K_N\ge0$ and $\frac12\int_{-1}^1K_N=1$.
The finite geometric identity implies, for $0<|t|\le1$,
\[
 K_N(t)\le\frac1{N S(t/2)^2}\le\frac1{Nt^2}.
\]
For the last inequality, the midpoint identity makes $S$ concave on
$[0,1]$, where $S\ge0$: its centered second difference is
$2S(x)(C(h)-1)\le0$.  Rational Jensen inequalities on $[0,1/2]$ give
$S(t/2)\ge t$ for $0\le t\le1$.

The kernel computation
\[
 \sigma_N f(x)=\frac12\int_{-1}^1K_N(t)f(x-t)\,dt
\]
is also the finite trigonometric sum
\[
 \sum_{|k|<N}\left(1-\frac{|k|}{N}\right)\widehat f(k)e_k(x).
\]
This follows by expanding the finite kernel and shifting the integral;
periodicity makes any interval of length $2$ equivalent to $[-1,1]$.
The shift is justified by splitting at the finitely many period boundaries.

Let $|f|\le B$, and choose rational $0<\delta\le1$ so that periodic
arguments within $\delta$ have values within $\varepsilon$.  Split the
kernel integral at $|t|=\delta$.  Positivity and normalization bound the
near part by $\varepsilon$; the far part is at most
$2B/(N\delta^2)$.  Hence
\[
 \boxed{|\sigma_N f(x)-f(x)|\le\varepsilon+\frac{2B}{N\delta^2}.}
\]
This constructs a uniform trigonometric approximation with a terminating
precision schedule.  Computing the finitely many Fourier coefficients to
additional error $\eta/(2N)$ gives at most $\eta$ further output error.

\section{The periodic square}
Let $f(x)=x^2$ on $[-1,1]$, extended with period $2$.  Its endpoints match,
and it has periodic Lipschitz bound $2$, verified by splitting an increment
at the join if necessary.  Two integrations by parts give
\[
 \widehat f(0)=\frac13,\qquad
 \widehat f(k)=\frac{2(-1)^k}{\pi^2k^2}\quad(k\ne0).
\]
Here is the real calculation behind the coefficient: symmetry removes the
sine term, and with $a=\pi k$,
\[
 \int_0^1x^2\cos(ax)\,dx
 =-\frac2a\int_0^1x\sin(ax)\,dx
 =\frac{2(-1)^k}{a^2}.
\]
All boundary terms are explicit geometric sine and cosine values.

The series
\[
 F(x)=\frac13+\frac4{\pi^2}
          \sum_{k=1}^{\infty}\frac{(-1)^kC(kx)}{k^2}
\]
has a uniform tail at most $4/(\pi^2N)$ after term $N$, by the earlier
reciprocal-square tail.  It is therefore an independent continuous function
computation.  Its Fejer sums are the same finite sums as those of $f$:
termwise integration follows from the uniform tail, and finite
orthogonality gives the coefficients.  The kernel estimate shows that
these common sums converge to both $F$ and $f$.  Thus $F=f$.

\begin{theorem}[Basel identity]\label{thm:basel-geometric-overlap}
The reciprocal-square computation of Chapter~\ref{ch:infinite-series}
satisfies
\[
 \boxed{\sum_{k=1}^{\infty}\frac1{k^2}=\frac{\pi^2}{6}.}
\]
\end{theorem}
\begin{proof}
At the rational endpoint $x=1$, $C(k)=(-1)^k$ and $f(1)=1$.  The proved
series identity gives $1=1/3+(4/\pi^2)\sum_{k\ge1}1/k^2$.
\end{proof}

The chain of comparisons is now complete: rational circle areas define
$\pi$; convex secants compute sine and cosine derivatives; the FTC computes
Fourier coefficients; finite positive kernels identify a uniformly
convergent series; evaluating it identifies an independently constructed
number.  No Hilbert-space completeness or unrestricted Fourier convergence
theorem was needed.
